\documentclass[10pt,a4paper]{amsart}
\usepackage{amsmath,amssymb,amsthm,graphicx}
% Retained minimal preamble for the Stacks Project chapter more-algebra.tex (More on Algebra).
% The chapter file pinned-original-source.tex carries only \input{preamble} on its first line, so no
% preamble text lives inside the chapter itself.  The Stacks Project uses one single shared preamble for
% every chapter; the snapshot stages it identically as sources/stacks-chow/preamble.tex and
% sources/stacks-etale-cohomology/preamble.tex.  The block below is copied mechanically, line for line,
% from sources/stacks-chow/preamble.tex lines 14-18 (the Xy-pic setup) and 154-217 (the native theorem
% declarations together with the mathematical macros).  Nothing was added, removed or redefined; the
% wrapper class and the amsmath/amssymb/amsthm/graphicx base packages are supplied by the assembler.
% Chapter: more-algebra.tex -- More on Algebra.
% For commutative diagrams we use Xy-pic
\usepackage[all]{xy}

% We use 2cell for 2-commutative diagrams.
\xyoption{2cell}

% Theorem environments.
%
\theoremstyle{plain}
\newtheorem{theorem}[subsection]{Theorem}
\newtheorem{proposition}[subsection]{Proposition}
\newtheorem{lemma}[subsection]{Lemma}

\theoremstyle{definition}
\newtheorem{definition}[subsection]{Definition}
\newtheorem{example}[subsection]{Example}
\newtheorem{exercise}[subsection]{Exercise}
\newtheorem{situation}[subsection]{Situation}

\theoremstyle{remark}
\newtheorem{remark}[subsection]{Remark}
\newtheorem{remarks}[subsection]{Remarks}

\numberwithin{equation}{subsection}

% Macros
%
\def\lim{\mathop{\mathrm{lim}}\nolimits}
\def\colim{\mathop{\mathrm{colim}}\nolimits}
\def\Spec{\mathop{\mathrm{Spec}}}
\def\Hom{\mathop{\mathrm{Hom}}\nolimits}
\def\Ext{\mathop{\mathrm{Ext}}\nolimits}
\def\SheafHom{\mathop{\mathcal{H}\!\mathit{om}}\nolimits}
\def\SheafExt{\mathop{\mathcal{E}\!\mathit{xt}}\nolimits}
\def\Sch{\mathit{Sch}}
\def\Mor{\mathop{\mathrm{Mor}}\nolimits}
\def\Ob{\mathop{\mathrm{Ob}}\nolimits}
\def\Sh{\mathop{\mathit{Sh}}\nolimits}
\def\NL{\mathop{N\!L}\nolimits}
\def\CH{\mathop{\mathrm{CH}}\nolimits}
\def\proetale{{pro\text{-}\acute{e}tale}}
\def\etale{{\acute{e}tale}}
\def\QCoh{\mathit{QCoh}}
\def\Ker{\mathop{\mathrm{Ker}}}
\def\Im{\mathop{\mathrm{Im}}}
\def\Coker{\mathop{\mathrm{Coker}}}
\def\Coim{\mathop{\mathrm{Coim}}}

% Boxtimes
%
\DeclareMathSymbol{\boxtimes}{\mathbin}{AMSa}{"02}

%
% Macros for moduli stacks/spaces
%
\def\QCohstack{\mathcal{QC}\!\mathit{oh}}
\def\Cohstack{\mathcal{C}\!\mathit{oh}}
\def\Spacesstack{\mathcal{S}\!\mathit{paces}}
\def\Quotfunctor{\mathrm{Quot}}
\def\Hilbfunctor{\mathrm{Hilb}}
\def\Curvesstack{\mathcal{C}\!\mathit{urves}}
\def\Polarizedstack{\mathcal{P}\!\mathit{olarized}}
\def\Complexesstack{\mathcal{C}\!\mathit{omplexes}}
% \Pic is the operator that assigns to X its picard group, usage \Pic(X)
% \Picardstack_{X/B} denotes the Picard stack of X over B
% \Picardfunctor_{X/B} denotes the Picard functor of X over B
\def\Pic{\mathop{\mathrm{Pic}}\nolimits}
\def\Picardstack{\mathcal{P}\!\mathit{ic}}
\def\Picardfunctor{\mathrm{Pic}}
\def\Deformationcategory{\mathcal{D}\!\mathit{ef}}

% --- appended 2026-10-12 (Chief): the project own definitions of the reference/slogan/history
% display environments, copied mechanically from the shared Stacks preamble lines 8-12
% (\usepackage{verbatim} plus the three \newenvironment lines, which the chapter preamble keeps
% verbatim from the project preamble). Nothing added, removed or redefined: these environments
% carry no mathematical content and the project itself typesets them as comments. Without these
% three lines a counted statement whose own text uses them fails with Environment slogan undefined.
\usepackage{verbatim}
\newenvironment{reference}{\comment}{\endcomment}
\newenvironment{slogan}{\comment}{\endcomment}
\newenvironment{history}{\comment}{\endcomment}
\usepackage{cleveref}
\crefname{originallemma}{original lemma}{original lemmas}
\Crefname{originallemma}{Original lemma}{Original lemmas}
\makeatletter
\newlabel{lemma-cut-complex-in-two}{{068U}{source}{}{sourcepointer-1}{}}
\newlabel{lemma-cut-complex-in-two@cref}{{[originallemma][0][]068U}{[1][0][]source}{}{}{}}
\makeatother
\crefname{originallemma}{original lemma}{original lemmas}
\Crefname{originallemma}{Original lemma}{Original lemmas}
\makeatletter
\newlabel{lemma-lift-perfect-from-residue-field}{{0BCD}{source}{}{sourcepointer-2}{}}
\newlabel{lemma-lift-perfect-from-residue-field@cref}{{[originallemma][0][]0BCD}{[1][0][]source}{}{}{}}
\makeatother
\title{A bounded below pseudo-coherent complex whose fibre vanishes in low degrees becomes perfect after localizing}
\author{The Stacks Project Authors}
\date{stacks-project:more-algebra.tex:local-snapshot-20261012; GNU FDL 1.2 or later}
\begin{document}
\maketitle
\setcounter{subsection}{1}
\noindent\textit{Editorial note.} All delivered mathematical text is the author's own, copied from the pinned Stacks Project chapter file more-algebra.tex; no statement or proof step was rewritten, reordered or removed. The counted span is the author environment and its own proof environment. References inside the span are handed to the assembler as navigation stubs that carry the canonical tag of the cited tag only; whether each of them is load-bearing is recorded on the stub itself.\par\medskip
\noindent\textit{Allowed standard prerequisites (editorial).}\begin{itemize}
\item After replacing the ring by a localization at a suitable element, a pseudo-coherent complex splits as the direct sum of its truncation below a degree and its truncation above that degree, the latter being perfect. invoked at proof lines 20020-20024. Produces the splitting of the complex and the perfectness of the upper truncation.
\item A perfect complex whose fibre over a prime is a finite free complex concentrated in degrees where the dimensions are prescribed can be represented over a suitable localization by the complex with those free modules. invoked at proof lines 20027-20029. Turns the perfect object into the displayed bounded complex of free modules with the prescribed ranks.
\end{itemize}\medskip
\ifdefined\hypertarget\hypertarget{sourcepointer-1}{}\fi
\noindent\textit{Source reference.} Original lemma 068U: sources/stacks-more-algebra/pinned-original-source.tex, source lines 19867--19867. Invoked at proof lines 20020-20024 to split the complex and to get the perfectness of the upper truncation; the vanishing of the lower truncation and the conclusion are inside the counted span.\par\medskip
\ifdefined\hypertarget\hypertarget{sourcepointer-2}{}\fi
\noindent\textit{Source reference.} Original lemma 0BCD: sources/stacks-more-algebra/pinned-original-source.tex, source lines 19594--19594. Used at proof lines 20027-20029 to produce the displayed free complex representing the object; the reduction to the perfect case that makes it applicable is inside the counted span.\par\medskip
\begin{lemma}
\label{lemma-lift-bounded-pseudo-coherent-to-perfect}
Let $R$ be a ring and let $\mathfrak p \subset R$ be a prime.
Let $K$ be pseudo-coherent and bounded below. Set
$d_i = \dim_{\kappa(\mathfrak p)}
H^i(K \otimes_R^\mathbf{L} \kappa(\mathfrak p))$.
If there exists an $a \in \mathbf{Z}$ such that $d_i = 0$ for $i < a$,
then there exists an $f \in R$, $f \not \in \mathfrak p$ and a
complex
$$
\ldots \to 0 \to R_f^{\oplus d_a} \to R_f^{\oplus d_{a + 1}} \to
\ldots \to
R_f^{\oplus d_{b - 1}} \to
R_f^{\oplus d_b} \to 0
\to \ldots
$$
representing $K \otimes_R^\mathbf{L} R_f$ in $D(R_f)$.
In particular $K \otimes_R^\mathbf{L} R_f$ is perfect.
\end{lemma}

\begin{proof}
After decreasing $a$ we may assume that also $H^i(K^\bullet) = 0$ for
$i < a$. By Lemma \ref{lemma-cut-complex-in-two} after replacing $R$ by $R_f$
for some $f \in R$, $f \not \in \mathfrak p$
we can write $K^\bullet = \tau_{\leq a - 1}K^\bullet \oplus
\tau_{\geq a}K^\bullet$ in $D(R)$ with $\tau_{\geq a}K^\bullet$
perfect. Since $H^i(K^\bullet) = 0$ for $i < a$ we see that
$\tau_{\leq a - 1}K^\bullet = 0$ in $D(R)$. Hence $K^\bullet$
is perfect. Then we can conclude using
Lemma \ref{lemma-lift-perfect-from-residue-field}.
\end{proof}

\end{document}
